Three fixed sincere rankings $ABC$, $BCA$, $CAB$ give a majority cycle: $A$ beats $B$, $B$ beats $C$, and $C$ beats $A$, each 2--1. A vote chooses the higher of the two currently offered options on each card. Every individual ranking is consistent; the pairwise majority comparisons need not form a consistent ranking.

If two options vote first and their winner meets the waiting option, this same profile lets any option win by waiting. No card or ballot rule changes. Conversely, any option that beats both others wins every such three-option knockout. This finite example does not prove Arrow's theorem and makes no claim about fairness or real elections.

\textbf{Five fixed distance voters on one common line have a strict winner at the third favorite.} Order favorites $p_1\leq\cdots\leq p_5$ with multiplicity and assume $p_3$ is available. Against any $q<p_3$, voters 3,4,5 strictly prefer $p_3$; against any $q>p_3$, voters 1,2,3 strictly prefer it. Equal-distance ties among the other voters cannot defeat these three votes. This proves a unique strict pairwise winner against every other available location. A favorite at zero is a voter. Repeated favorites are separate voters, not one position. The middle favorite need not be the mean or the midpoint of the printed line.

For four favorites 1,3,7,9 on the line, every point in $[3,7]$ never loses a comparison, but no point beats every other. For the integer board these weak winners are 3,4,5,6,7. The even-voter result is different from saying every vote ties. With three distinct candidate positions on a line, the middle candidate cannot be any distance voter's strict worst choice. Since the cyclic profile puts each candidate last once, no ordering of candidate positions realizes all three cards by nearest-distance preferences.
